\section{Autoregressive Factorization}
\paragraph{Purpose}
A language model is a probability distribution over sequences of tokens
$\bm{x}=x_{1:n}$. By the chain rule of probability, \uB{with no approximation}:
\begin{center}
    \frB{$\mathbb{P}\left(x_{1:n}\right) = \prd{t=1}{n}
    \mathbb{P}\left(x_{t}\,|\,x_{1:t-1}\right)$}
\end{center}
so it is enough to model the distribution of the \tB{next token given the previous ones}.
The network outputs a vector of logits $\bm{z}_{t}\in\mathbb{R}^{V}$ and
$\mathbb{P}_{\bm{\theta}}\left(x_{t}=k\,|\,x_{<t}\right)=\mathrm{softmax}(\bm{z}_{t})_{k}$.

\paragraph{Relationship with Markov models}
An $n$-gram model assumes that the next token depends only on the $n-1$ previous ones, that is
a Markov chain of order $n-1$, and estimates the probabilities by counting. The table has
$V^{n}$ entries: this is the \tB{curse of dimensionality}, most contexts are never seen.
A neural language model conditions on the whole context and shares its parameters across
positions, and similar contexts get similar representations, which is how it generalizes.

\section{Training Objective}
\paragraph{Theory}
For a corpus of $N$ token positions, the maximum likelihood estimate minimizes the
\tB{cross-entropy} (the average negative log-likelihood of the next token):
\begin{center}
    \frB{$\mathcal{L}(\bm{\theta}) = -\dfrac{1}{N}\su{t=1}{N}\log
    \mathbb{P}_{\bm{\theta}}\left(x_{t}\,|\,x_{<t}\right)$}
\end{center}
This is the multi-class classification loss of the feed-forward networks (softmax and
cross-entropy) applied at every position, with a vocabulary of $V$ classes.
\begin{itemize}
    \item it is \textbf{self-supervised}: the labels are the text itself, shifted by one
        position, so no annotation is needed;
    \item with the causal mask, the inputs $x_{1:n-1}$ give the $n-1$ predictions
        of $x_{2:n}$ in \uB{one parallel pass} (\emph{teacher forcing}: the true
        prefix is fed, not the model's own predictions);
    \item for the data distribution $p$: $H(p,q_{\bm{\theta}}) = H(p) + KL\left(p\,\|\,
        q_{\bm{\theta}}\right)$, so minimizing the cross-entropy minimizes the KL
        divergence to the data, and its minimum is the entropy of the text, which is not zero.
\end{itemize}

\section{Perplexity}
\paragraph{Theory}
The standard evaluation metric is the exponential of the cross-entropy (in nats):
\begin{center}
    \fr{$PPL = \exp\left(\mathcal{L}\right) = \left(\prd{t=1}{N}
    \mathbb{P}_{\bm{\theta}}\left(x_{t}\,|\,x_{<t}\right)\right)^{-1/N}$}
\end{center}
Interpretation: the model is as uncertain as if it chose uniformly among $PPL$ tokens at
each step. A uniform model over the vocabulary has $PPL=V$ and a perfect one has $PPL=1$.

\paragraph{Weaknesses}
\begin{itemize}
    \item it depends on the tokenizer: two models with different vocabularies cannot be
        compared directly (use \emph{bits per byte} instead);
    \item it depends on the domain of the evaluation text;
    \item a low perplexity does not mean the answers are helpful or true.
\end{itemize}
